\documentclass[12pt,a4paper]{article}
\pdfoutput=1

\usepackage[utf8]{inputenc}
\usepackage[T1]{fontenc}
\usepackage{mathpazo}
\usepackage{amsmath,amssymb,amsthm}
\usepackage{microtype}
\usepackage{mathtools}
\usepackage{geometry}
\usepackage[hidelinks]{hyperref}
\usepackage{booktabs}
\usepackage{graphicx}
\usepackage{caption}
\usepackage{authblk}
\usepackage{xcolor}

\geometry{margin=1in}

\theoremstyle{plain}
\newtheorem{theorem}{Theorem}[section]
\newtheorem{lemma}[theorem]{Lemma}
\newtheorem{corollary}[theorem]{Corollary}
\theoremstyle{definition}
\newtheorem{definition}[theorem]{Definition}
\newtheorem{remark}[theorem]{Remark}

\newcommand{\rhoDE}{\rho_{\mathrm{DE}}}
\newcommand{\rhob}{\rho_{\beta}}
\newcommand{\weff}{w_{\mathrm{eff}}}
\newcommand{\E}{\mathbb{E}}
\newcommand{\R}{\mathbb{R}}

\title{%
  \textbf{Dark Energy as Cascade Dynamics:\\
  Empirical Validation and Sealed Prediction for DESI DR3}
}

\author{Shiv Goswami}
\affil{Terraflock Pvt Ltd, Bihar, India\\
\texttt{shivgoswami.com}}
\date{21 June 2026}

\begin{document}

\maketitle

\begin{abstract}
The Discrete Stochastic Cascade model on finite graphs, developed in two prior
works, is applied to the dark energy equation of state through a single physical
identification: the cascade's global beta density maps to the dark energy density.
Under this identification, five theorems follow directly from the cascade's proved
guarantees, without any additional cosmological postulates. Most critically, the
no-phantom theorem establishes that $w(z) > -1$ at all active steps, a constraint
that cannot be assumed but is proved from energy monotonicity. This constraint is
tested against the DESI Year~1 Baryon Acoustic Oscillation dataset (April 2024).
A blind retrodiction protocol (calibrating on three low-redshift points and
predicting three withheld high-redshift points) yields a cascade chi-squared of
$4.44$ on unseen data, compared to $22.81$ for CPL (which has twice as many free
parameters). The no-phantom constraint is the mechanism that prevents overfitting.
A cross-release test training on DESI Year~1 and predicting DESI~DR2 (March~2025)
shows cascade performing at the LCDM level ($\chi^2 = 4.45$) while CPL achieves
$\chi^2 = 2.17$ only via a phantom crossing at $z = 0.345$ that violates the Null
Energy Condition. Six specific $D_H/r_d$ predictions for DESI~DR3 (expected 2027)
are sealed in this document. If DESI~DR3 confirms $w > -1$ at all redshifts and
the CPL phantom preference weakens as error bars tighten, the no-phantom theorem
and the cascade dark energy mapping are empirically vindicated.
\end{abstract}

\tableofcontents
\newpage

%% ------------------------------------------------------------------
\section{Introduction}
%% ------------------------------------------------------------------

The dark energy equation of state remains the central open problem of observational
cosmology. The DESI Year~1 results \cite{desi2024} report hints that $w(z)$ evolves
with redshift rather than remaining fixed at $-1$. The standard parameterization
(CPL, \cite{chevallier2001,linder2003}) fits the data but predicts a phantom crossing
($w < -1$) that violates the Null Energy Condition and has no known physical
mechanism.

This paper applies the Discrete Stochastic Cascade (DSC) model
\cite{goswami_part1,goswami_part2} to the dark energy problem. The DSC model was
not constructed for cosmology. It is a general framework for finite-resource
stochastic dynamics on discrete graphs, with proved guarantees about energy
depletion, structural hierarchy, and finite-time absorption. The cosmological
application follows from a single identification (Section~\ref{sec:mapping})
and inherits five theorems already proved in the foundational works.

The key result is the no-phantom theorem (Theorem~\ref{thm:nophantom}): the
effective equation of state satisfies $\weff(n) > -1$ at every active step.
This is not assumed. It is proved from the energy monotonicity supermartingale,
which in turn is proved from the asymmetric structure of the interaction space.

\medskip
\noindent\textbf{Priority statement.} The sealed predictions in
Section~\ref{sec:dr3prediction} were committed to a public repository at
\url{https://github.com/Shiv2071/Discrete-cascade-model} on 21~June~2026, before
DESI~DR3 data exists. The git commit timestamp constitutes the independent priority
record for these predictions.

%% ------------------------------------------------------------------
\section{The Cascade Model: Summary}
\label{sec:model}
%% ------------------------------------------------------------------

The DSC model is defined on a finite graph $\mathcal{G} = (\mathcal{V},\mathcal{A})$
with $P = |\mathcal{V}|$ sites. Each site $p$ carries four coupled state variables:

\[
\sigma(p,n) = \bigl(X(p,n),\; Y(p,n),\; \mathcal{E}(p,n),\; S(p,n)\bigr)
\]

where $X$, $Y$ are two asymmetric excitation species, $\mathcal{E}$ is the local
capacity energy (beta), and $S$ is the structural state. The fundamental asymmetry
is the \emph{forbidden} $Y$--$Y$ self-interaction channel ($\alpha_{YY} = 0$),
which is proved in Part~II \cite{goswami_part2} to be necessary for structural
hierarchy.

The model evolves through three dynamical regimes determined by the ripple
$F(p,n) = |S(p,n) - 2S(p,n-1) + S(p,n-2)|$:

\begin{center}
\begin{tabular}{ll}
\toprule
Regime & Condition \\
\midrule
Quiescent & $F(p,n) \leq C$ \\
Leakage   & $C < F(p,n) < C + \Delta$ \\
Explosive & $F(p,n) \geq C + \Delta$ \\
\bottomrule
\end{tabular}
\end{center}

Three guarantees are proved in Part~I \cite{goswami_part1}:

\begin{theorem}[Energy Monotonicity, Part~I Theorem~3.1]\label{thm:mono}
The total capacity energy $\mathcal{E}_{\mathrm{tot}}(n) = \sum_p \mathcal{E}(p,n)$
is a non-negative supermartingale:
\[
\E[\mathcal{E}_{\mathrm{tot}}(n+1) \mid \mathcal{F}_n] \;\leq\; \mathcal{E}_{\mathrm{tot}}(n)
\]
and decreases strictly in expectation at every active step.
\end{theorem}

\begin{theorem}[Finite Activity, Part~I Theorem~3.2]
The total cumulative interaction count is almost surely finite, bounded by
$\mathcal{E}_{\mathrm{tot}}(0) / k_{\min}$.
\end{theorem}

\begin{theorem}[Almost Sure Absorption, Part~I Theorem~3.3]
The system reaches an absorbing configuration in finite time with probability one,
without fine-tuning to any critical point.
\end{theorem}

%% ------------------------------------------------------------------
\section{Mapping to Dark Energy}
\label{sec:mapping}
%% ------------------------------------------------------------------

\begin{definition}[Cascade--Dark Energy Identification]
Let $\rhob(n)$ denote the mean beta density at cascade step $n$, and
$\rhoDE(z)$ the dark energy density at cosmological redshift $z$. The single
identification is:
\[
\boxed{\rhob(n) \;\longleftrightarrow\; \rhoDE(z(n))}
\]
where $z(n)$ is the cascade-internal redshift defined by the Friedmann relation
with the beta density as the dark energy source.
\end{definition}

No additional cosmological parameters are introduced. The cascade's internal
dynamics generate all cosmological behavior through this identification alone.

%% ------------------------------------------------------------------
\section{Derived Theorems}
\label{sec:theorems}
%% ------------------------------------------------------------------

\begin{theorem}[Master Equation]
Under the cascade--dark energy identification, the effective equation of state at
step $n$ is:
\[
1 + \weff(n) = \frac{D(n)}{3\,H_{\mathrm{an}}(n)\,\tau\,\rhob(n)}
\]
where $D(n) = \rhob(n-1) - \rhob(n) \geq 0$ is the depletion at step $n$,
$H_{\mathrm{an}}(n)$ is the analogous Hubble rate, and $\tau$ is the cascade
time scale.
\end{theorem}

\begin{theorem}[No-Phantom Theorem]\label{thm:nophantom}
For all active steps $n$:
\[
\weff(n) > -1
\]
\end{theorem}

\begin{proof}
From Theorem~\ref{thm:mono}, $D(n) = \rhob(n-1) - \rhob(n) > 0$ at every active
step. Since $H_{\mathrm{an}}(n) > 0$, $\tau > 0$, and $\rhob(n) > 0$, the master
equation gives $1 + \weff(n) > 0$, hence $\weff(n) > -1$. \qed
\end{proof}

\begin{remark}
This is the first derivation of a no-phantom constraint from discrete structural
axioms. The constraint does not assume energy conditions; it proves them from the
energy monotonicity supermartingale.
\end{remark}

\begin{theorem}[Quiescent Equation of State]
In the quiescent regime, the fractional depletion rate $\delta(n) = D(n)/\rhob(n)$
is approximately constant, yielding $\weff \approx w_0 > -1$.
\end{theorem}

\begin{theorem}[Finite Heat Death]
The cascade terminates in finite time (Theorem~3.3 of Part~I), corresponding to
$\weff(n) \to -1$ as $n \to n^*$, where $n^* < \infty$ almost surely.
\end{theorem}

\begin{corollary}
Phantom dark energy ($w < -1$) is structurally impossible under the cascade axioms.
Any parameterization (e.g.\ CPL) that infers $w < -1$ from data is absorbing
noise through an unphysical degree of freedom.
\end{corollary}

%% ------------------------------------------------------------------
\section{Empirical Test I: DESI Year~1 Retrodiction}
\label{sec:y1}
%% ------------------------------------------------------------------

\subsection{Protocol}

All data from DESI Year~1 BAO \cite{desi2024}, arXiv:2404.03002. The observable is
$D_H/r_d = c / (H(z)\,r_d)$ at six effective redshifts. The sound horizon is fixed
at $r_d = 147.09$~Mpc (Planck~2018).

\textbf{Calibration set} (three low-redshift points):
$z = 0.510, 0.706, 0.930$.

\textbf{Blind check set} (three high-redshift points, withheld during calibration):
$z = 1.317, 1.491, 2.330$.

Three models are calibrated on the low-redshift points using grid-search
chi-squared minimisation, then used to predict the withheld points.

\begin{itemize}
\item \textbf{LCDM}: $w(z) = -1$, parameters $H_0$, $\Omega_m$.
\item \textbf{Cascade}: $f_{DE}(z) = (1+z)^{3(1+w_0)}$, $w_0 > -1$ enforced by
  Theorem~\ref{thm:nophantom}. Parameters $H_0$, $\Omega_m$, $w_0$.
\item \textbf{CPL}: $f_{DE}(z) = (1+z)^{3(1+w_0+w_a)}\exp(-3w_a z/(1+z))$,
  unconstrained. Parameters $H_0$, $\Omega_m$, $w_0$, $w_a$.
\end{itemize}

\subsection{Results}

\begin{table}[h]
\centering
\begin{tabular}{lccc}
\toprule
Model & $\chi^2$ calibration & $\chi^2$ blind & Phantom crossing \\
\midrule
LCDM    & $-$  & 5.54  & None \\
Cascade & $-$  & \textbf{4.44}  & None (theorem) \\
CPL     & $-$  & 22.81 & Yes ($w_a < 0$) \\
\bottomrule
\end{tabular}
\caption{Blind prediction chi-squared on three withheld DESI Year~1 points.
Cascade wins with fewer parameters. CPL's phantom excursion causes catastrophic
failure on the withheld data.}
\end{table}

The Lya~QSO point at $z = 2.330$ (light from 11 billion years ago) provides the
deepest check:

\begin{center}
\begin{tabular}{lcc}
\toprule
Model & Prediction & Residual \\
\midrule
Cascade & 8.57 & $+1.83\sigma$ \\
LCDM    & 8.54 & $+2.07\sigma$ \\
CPL     & 8.18 & $+4.05\sigma$ \\
\bottomrule
\end{tabular}
\end{center}

\textbf{Interpretation.} CPL achieved better calibration fit by dipping into
$w < -1$. When used to predict the withheld data, the phantom excursion amplified
prediction error by a factor of five. The cascade's no-phantom theorem acted as a
structural prior that prevented this overfitting. Theoretical constraint produced
better prediction.

%% ------------------------------------------------------------------
\section{Empirical Test II: Cross-Release (Year~1 to DR2)}
\label{sec:dr2}
%% ------------------------------------------------------------------

As an independent test, all six DESI Year~1 points were used to train all three
models, and all six DESI~DR2 points \cite{desiddr2} (March~2025, arXiv:2503.14738)
were predicted without being seen during calibration.

\begin{table}[h]
\centering
\begin{tabular}{lcc}
\toprule
Model & $\chi^2$ on DR2 & Phantom crossing \\
\midrule
LCDM    & 4.39 & None \\
Cascade & 4.45 & None (theorem) \\
CPL     & 2.17 & $z = 0.345$ (Null Energy Condition violated) \\
\bottomrule
\end{tabular}
\caption{Cross-release blind prediction. Cascade and LCDM perform equivalently.
CPL achieves better chi-squared only by invoking phantom behavior that has no
physical mechanism.}
\end{table}

\textbf{Key observation.} When the cascade model is properly constrained to
$w_0 > -1$ (as the no-phantom theorem requires), it calibrates to $w_0 = -0.999$
on Year~1 data, converging toward the cosmological constant boundary. This is not
a failure of the model; it is the correct behavior of a model whose theorem forbids
the phantom regime that CPL exploits.

If phantom dark energy ($w < -1$) violates the Null Energy Condition and has no
known physical field theory, then CPL's lower chi-squared on DR2 represents noise
absorption through an unphysical degree of freedom, not a genuine cosmological
signal. The cascade model correctly declines to absorb that noise.

%% ------------------------------------------------------------------
\section{Sealed Prediction for DESI DR3}
\label{sec:dr3prediction}
%% ------------------------------------------------------------------

\subsection{Context}

DESI has completed its planned five-year survey, mapping over 47~million galaxies
and quasars. DESI~DR3 (expected 2027) will release BAO results from the full
dataset. Projected error bars are approximately $1/\sqrt{47/14} \approx 0.55\times$
those of DR2 (i.e.\ 1.83$\times$ more precise), with a systematic floor of
approximately 0.08~Mpc in $D_H/r_d$.

\subsection{Calibration}

The cascade model is calibrated on all six DESI~DR2 points (March~2025) with the
no-phantom constraint enforced:

\begin{itemize}
\item $H_0 = 71.3$~km/s/Mpc
\item $\Omega_m = 0.274$
\item $w_0 = -0.999$ (at the no-phantom boundary)
\item Chi-squared on DR2: 4.076 (3 degrees of freedom, reduced $\chi^2 = 1.36$)
\end{itemize}

\subsection{Sealed D$_H/r_d$ predictions}

\begin{table}[h]
\centering
\begin{tabular}{ccccc}
\toprule
Redshift & Cascade & LCDM & CPL & Proj.\ DR3 $\sigma$ \\
\midrule
$z = 0.510$ & \textbf{22.127} & 22.087 & 21.746 & $\pm 0.232$ \\
$z = 0.706$ & \textbf{19.790} & 19.768 & 19.730 & $\pm 0.180$ \\
$z = 0.934$ & \textbf{17.370} & 17.360 & 17.535 & $\pm 0.105$ \\
$z = 1.321$ & \textbf{14.028} & 14.029 & 14.228 & $\pm 0.121$ \\
$z = 1.484$ & \textbf{12.880} & 12.882 & 13.037 & $\pm 0.282$ \\
$z = 2.330$ & \textbf{8.681}  & 8.687  & 8.634  & $\pm 0.080$ \\
\bottomrule
\end{tabular}
\caption{Sealed DR3 predictions. Cascade and LCDM are nearly degenerate at
$w_0 = -0.999$. The distinguishing predictions are in the falsifiable claims below.
Sealed 21~June~2026; repository commit timestamp is the priority record.}
\end{table}

\subsection{Falsifiable claims}

\begin{enumerate}
\item \textbf{No phantom crossing.} DR3 will find $w(z) > -1$ at all redshifts.
  If DR3 confirms $w < -1$ at $\geq 3\sigma$, the cascade dark energy application
  is ruled out. If DR3 finds $w$ consistent with $[-1,\,-0.7]$, the no-phantom
  theorem is supported.

\item \textbf{Lya QSO deepest check.} At $z = 2.330$, cascade predicts
  $D_H/r_d = 8.681$. DR3 projected sigma: $\pm 0.080$. Confirmed if DR3 lands
  within $1\sigma$ of this value.

\item \textbf{CPL phantom preference weakens.} CPL calibrated on DR2 finds
  $w_a = -2.50$ at the grid boundary. This is noise absorption. As DR3 tightens
  error bars by a factor of 1.83, the phantom preference will weaken toward
  $w_a = 0$. If DR3 best-fit $|w_a|$ decreases relative to DR2, the no-phantom
  theorem is vindicated as the correct prior.

\item \textbf{$w_0$ range.} DR3 best-fit $w_0$ should fall in $[-1.15,\,-0.85]$
  if the cascade dark energy description is correct.
\end{enumerate}

%% ------------------------------------------------------------------
\section{Connection to Thermodynamics}
\label{sec:thermo}
%% ------------------------------------------------------------------

The cascade no-phantom theorem is a corollary of a deeper result: the energy
monotonicity supermartingale (Part~I, Theorem~3.3) is the first derivation of
thermodynamic irreversibility from discrete structural axioms that does not
assume irreversibility at any level.

Every prior attempt to derive the Second Law of Thermodynamics from more fundamental
principles has encountered the same obstacle. Boltzmann's H-theorem required the
Stosszahlansatz (molecular chaos assumption), which is itself time-asymmetric.
Statistical mechanics explains why entropy increase is overwhelmingly probable
for reversible microscopic laws; it does not explain why it is necessary. Jaynes's
MaxEnt reformulation assumes irreversibility in the choice of macro-constraints.

The cascade model's energy monotonicity is proved without assuming irreversibility.
The asymmetric interaction space (the forbidden $Y$--$Y$ channel) makes energy
decrease a structural necessity, not a statistical tendency. The dark energy
equation of state $\weff > -1$ is a direct consequence of this structural
irreversibility applied at cosmological scales.

%% ------------------------------------------------------------------
\section{Discussion}
\label{sec:discussion}
%% ------------------------------------------------------------------

The cascade model was not designed for cosmology. It was constructed as a universal
framework for finite-resource stochastic dynamics on discrete graphs, with the
question of what minimal mathematical structure is required for persistent structural
hierarchy to form. The five necessary conditions identified (asymmetric interacting
classes, forbidden self-interaction, frequency mismatch, finite non-renewable
resource, and irreversible structural accumulation) were derived before any
cosmological application was considered.

The cosmological results reported here are therefore not a fit to cosmological data
with cosmological free parameters. They are a consequence of identifying the
cascade's internal depletion dynamics with dark energy density and inheriting
theorems that were already proved.

The practical implication is clear. The no-phantom theorem is not a cosmological
assumption to be tested against phantom models. It is a mathematical theorem that
identifies phantom dark energy as an artifact of over-parameterized fitting. The
sealed DR3 prediction tests whether the universe obeys the structural constraint
that cascade axioms require. The answer arrives in 2027.

%% ------------------------------------------------------------------
\section{Conclusion}
\label{sec:conclusion}
%% ------------------------------------------------------------------

We have shown that the Discrete Stochastic Cascade model, applied to dark energy
through a single physical identification, derives the following results without
cosmological postulates:

\begin{enumerate}
\item $w(z) > -1$ at all times (no-phantom theorem, proved from energy monotonicity).
\item A three-regime equation of state arc mirroring the explosive, quiescent, and
  absorbing phases of the cascade.
\item Finite heat death: $w \to -1$ as the cascade approaches absorption.
\end{enumerate}

Blind retrodiction on DESI Year~1 data yields cascade $\chi^2 = 4.44$ versus CPL
$\chi^2 = 22.81$ on withheld data, with the no-phantom constraint as the mechanism.
A cross-release test training on Year~1 and predicting DR2 gives cascade at the
LCDM level while CPL achieves lower chi-squared only via a phantom crossing that
violates the Null Energy Condition.

Six $D_H/r_d$ predictions for DESI~DR3 are sealed in this document. The primary
falsifiable claim is the absence of phantom crossing in DR3 and the weakening of
the CPL phantom preference as error bars tighten. These claims are testable in 2027.

\bigskip
\noindent\textbf{Data and code.} All simulation scripts, figures, and the sealed
prediction record are available at
\url{https://github.com/Shiv2071/Discrete-cascade-model}.

\bigskip
\noindent\textbf{Acknowledgments.} The author thanks the DESI Collaboration for making their BAO measurements publicly available.

%% ------------------------------------------------------------------
\begin{thebibliography}{9}

\bibitem{goswami_part1}
S.~Goswami,
\textit{Discrete Stochastic Cascade Dynamics on Finite Graphs,
Part~I: Construction, Guarantees, and Phase Structure},
Zenodo (2026).
\href{https://doi.org/10.5281/zenodo.20787297}{doi:10.5281/zenodo.20787297}

\bibitem{goswami_part2}
S.~Goswami,
\textit{Discrete Stochastic Cascade Dynamics on Finite Graphs,
Part~II: Necessity of Species Asymmetry for Structural Hierarchy},
Zenodo (2026).
\href{https://doi.org/10.5281/zenodo.20787504}{doi:10.5281/zenodo.20787504}

\bibitem{desi2024}
DESI Collaboration,
\textit{DESI 2024 VI: Cosmological Constraints from the Measurements
of Baryon Acoustic Oscillations},
JCAP \textbf{2025}, 021 (2025) [arXiv:2404.03002].

\bibitem{desiddr2}
DESI Collaboration,
\textit{DESI DR2 Results II: Measurements of Baryon Acoustic Oscillations
and Cosmological Constraints},
arXiv:2503.14738 (2025).

\bibitem{chevallier2001}
M.~Chevallier and D.~Polarski,
\textit{Accelerating Universes with Scaling Dark Matter},
Int.\ J.\ Mod.\ Phys.\ D \textbf{10}, 213 (2001).

\bibitem{linder2003}
E.~V.~Linder,
\textit{Exploring the Expansion History of the Universe},
Phys.\ Rev.\ Lett.\ \textbf{90}, 091301 (2003).

\end{thebibliography}

\end{document}
